\section{Foundation Models 与 Robotics 之间缺什么}

Foundation model 擅长 high-level reasoning，机器人系统却要在几十到几百毫秒内输出连续控制。二者之间不是一条自然存在的 API：模型输入输出的表示、训练数据的尺度、控制频率与安全约束都不同。本章把 lecture 的“为什么不能直接接上”完整展开。

\subsection{High-level reasoning 与 low-level control}

图中左侧的 PaLM-SayCan 负责从语言目标选择高层 skill，右侧 RT-1 负责把图像和指令转成低层动作。前者回答“下一步做什么”，后者回答“关节或末端执行器此刻怎样动”。高层计划可容忍秒级响应和离散步骤，低层 policy 必须处理连续误差与闭环修正。

\lecturefigure{slide-09-foundation-models-robotics-gap.jpg}{Foundation Models + Robotics：high-level reasoning 与 low-level control}{官方录像恢复幻灯片 V009；00:13:42}

一个低层 policy 可写成
$$
a_t\sim \pi_\theta(a_t\mid o_{\le t},g),
$$
其中 $a_t$ 是当前动作，$o_{\le t}$ 是截至时刻 $t$ 的视觉、proprioception 等 observation history，$g$ 是语言 goal，$\pi_\theta$ 是参数为 $\theta$ 的 policy。与一次性文本回答不同，这个分布会在每个控制周期重新调用，并改变下一步 observation。

\begin{importantbox}{First use：policy 与 trajectory}
\term{Policy（策略）}是从当前信息到动作分布的映射；\term{trajectory（轨迹）}是 $(o_0,a_0,o_1,a_1,\ldots)$ 的时序序列。Robot dataset 的基本单位常是 trajectory，而不是独立图文 pair。它更昂贵，因为每条数据需要真实或模拟 agent 与环境闭环交互。
\end{importantbox}

\subsection{从 modular stack 到 joint scaling}

时间线把 2014--2021 的模块化 perception/planning/control，与 2021--2023 的 RT-1、PaLM-E 和 RT-2 并列。趋势不是所有模块突然消失，而是更多表示和任务被放进共享参数空间，通过更大数据混合实现 model consolidation。下一页用 “combine large diverse offline datasets” 概括 scaling recipe，同时用 “Scale up and done?” 保留疑问。

\lecturefigure{slide-10-robotics-scaling-timeline.jpg}{机器人系统从模块化走向 model consolidation 的时间线}{官方录像恢复幻灯片 V010；00:15:06}

\lecturefigure{slide-11-scaling-recipe.jpg}{Scaling recipe：离线数据、多任务模型与 positive transfer}{官方录像恢复幻灯片 V011；00:15:36}

\begin{warningbox}{“Scale up” 不是缺省答案}
语言数据可从互联网被动采集，机器人数据需要 embodiment、环境、operator、reset 和安全成本。参数增加只有在 data mixture 提供相关 supervision、action representation 可学习、模型容量没有被错误任务占用时才可能产生 transfer。规模是实验变量，不是替代问题定义的口号。
\end{warningbox}

\subsection{Moravec's paradox：语言流畅不等于手指灵巧}

本节进一步解释为什么 high-level model success 不能直接外推到控制。Moravec's paradox 指出，对人类显得高阶的 symbolic reasoning 可能容易被计算机实现，而感知运动能力虽然主观上“无需思考”，却包含漫长进化和身体经验。机器人抓取中的毫米误差、接触切换、柔性物体和遮挡，往往比生成一段解释更难。

\lecturefigure{slide-12-moravec-paradox.jpg}{Moravec's paradox：低层感知运动可能比高层推理更难}{官方录像恢复幻灯片 V012；00:18:33}

\teachervoice{Fei Xia 用 Moravec's paradox 解释为什么 LLM 的成功没有自动解决机器人：会描述“怎样抓杯子”与能在摩擦变化、遮挡和接触冲击下稳定抓住杯子，是不同层级的能力。课程把这种落差作为 low-level embodied intelligence 的定义性难题。}

\subsection{两个根因：数据少与接口错位}

第一张 challenge slide 用数量级差异展示 robot trajectories 的稀缺；第二张则指出 LLM 输出文本 token，而机器人需要位置、旋转、gripper、速度或力矩。即使 LLM 拥有语义知识，如果训练中没有看到可对齐的动作监督，或者输出空间无法表达控制精度，也无法直接变成 policy。

\lecturefigure{slide-13-challenge-lack-of-data.jpg}{Challenge 1：robot data 远少于 image-text / text data}{官方录像恢复幻灯片 V013；00:20:21}

\lecturefigure{slide-14-challenge-llm-interface.jpg}{Challenge 2：LLM 没有天然的 low-level control interface}{官方录像恢复幻灯片 V014；00:21:30}

\lecturefigure{slide-15-agenda-two-parts.jpg}{两段路线：model consolidation 与 new interfaces}{官方录像恢复幻灯片 V015；00:23:27}

\begin{knowledgebox}{两条路线分别修什么}
PaLM-E / RT-2 路线主要修复\textbf{数据与表示共享}：让 web-scale semantic tasks 与 robot trajectories 在同一模型中共同训练。Language to Rewards 路线主要修复\textbf{接口错位}：语言模型不直接输出每个 actuator command，而输出一个可被控制器优化的 reward specification。
\end{knowledgebox}

\subsection{本章小结}

Foundation models 与 robot control 的缺口可归结为 data regime、action representation、control loop 和 safety constraint。Lecture 接下来先尝试把模型和数据整合进统一 VLA，再尝试保留传统 controller、只让 LLM 负责 reward interface。

